\section{可观测性与评估指标}

要让 Compound AI 在工程中稳定运行，需要把 trace、agent、solver 的运行状态映射到清晰的指标集合。

\subsection{关键指标面板}

\begin{itemize}
    \item \textbf{Trace fidelity}：trace 与 solver 输出 plan 的相似度（例如 Jaccard overlap）。
    \item \textbf{Judge reward curve}：agent judge 的 accuracy/proactivity/credibility 逐轮变化。
    \item \textbf{Constraint slip rate}：每份 plan 中 constraint violation 的出现频率。
    \item \textbf{Dialogue throughput}：APAC 问答完成一个 Json 需要的轮数。
\end{itemize}

\begin{importantbox}{指标背后的数据路径}
每个指标都应绑定到一个可追踪的数据源：trace fidelity 取自 search trace cache，judge reward curve 直接从 DPO judge 的 log stream 里采样，并且在 inference 侧还要记录 plan 的 \texttt{trace\_id}。Constraint slip rate 需要 solver 返回的 violations feed 到 observability db；Dialogue throughput 则要加上 latency instrumentation（下游 judge 会附带轮数、token count）。只有这样，dashboard 才能在秒级别内准确反映系统状态。
\end{importantbox}

\begin{knowledgebox}{指标的警戒线}
设置 guardrails：若 judge reward 连续三轮下降，就触发 new DPO fine-tune；若 trace fidelity < 0.7，就回退到 deterministic solver；若 constraint slip rate 上升，就加 \texttt{penalty\_update}。这些警戒线让运营可以快速捕捉模型 drift。
\end{knowledgebox}

\subsection{Trace-level 可靠性}

Trace 的可靠性可以通过覆盖率评估：每次 inference 都记录 trace ID，后端根据 trace metadata（e.g., depth, constraint violations）统计 coverage heatmap，确保训练集/测试集在每个 region 都有足够样本。

\begin{warningbox}{不要忽略 trace skew}
Long tail trace（极少数极长或有 large cost 的路径）很容易在训练中被忽略，但它们往往也是 failure 的来源。应建立 trace-skew dashboard，及时重采样这些样本。
\end{warningbox}

\subsection{本章小结}

具有 dashboard 级别的指标可以把抽象的 Compound AI pipeline 可观测化，一旦某个模块输出异常，就能沿着 trace → judge → solver 的路径迅速定位，并在 guardrail 触发时自动纠偏。

